\documentclass[11pt,a4paper]{article}

\usepackage[utf8]{inputenc}
\usepackage{amsmath}
\usepackage{graphicx}
\usepackage{hyperref}
\usepackage{booktabs}
\usepackage{geometry}
\usepackage{xcolor}

\geometry{margin=1in}

\title{\textbf{Project Report: ML-Guided Cut Selection for ASIC Technology Mapping in ABC}}
\author{Hemang Gautam}
\date{\today}

\begin{document}

\maketitle

\begin{abstract}
Technology mapping is a critical step in the ASIC design flow, heavily reliant on the generation and selection of structural ``cuts'' in an And-Inverter Graph (AIG). Selecting the right cuts minimizes the final Area-Delay Product (ADP). This project introduces a machine-learning-guided methodology to evaluate and select cuts within the ABC logic synthesis framework. Unlike traditional heuristic approaches, we utilize a Multi-Layer Perceptron (MLP) trained with a pairwise ranking loss to score cuts based on topological features. We integrated this model directly into ABC's C codebase via an ``Area Blending'' technique. This report details the methodology, feature engineering, deviations from existing literature (such as the SLAP framework), challenges faced during integration, and final empirical results demonstrating up to a 12.38\% ADP reduction on arithmetic benchmarks.
\end{abstract}

\section{Introduction}
In standard ASIC technology mapping (e.g., using ABC's \texttt{if} command), the mapper generates up to $C$ cuts per node (default $C=8$). It then evaluates these cuts using a localized, greedy heuristic (primarily Area Flow and Maximum Fanout Free Cone, or MFFC). While generating more cuts (e.g., $C=16$) increases the probability of finding a globally optimal mapping, it severely degrades CPU runtime. 

The goal of this project was to leverage machine learning to intelligently score and select cuts. If the ML model can accurately identify globally optimal cuts from a smaller generated pool, we can achieve high Quality of Results (QoR) without the runtime penalty of computing an exhaustive cut set. This report addresses three specific deliverables: analyzing mapping time versus cut capacity, detailing the ML model's architecture, and comparing the Area-Delay QoR against baseline ABC.

% ------------------------------------------------------------------
\section{Background \& Comparison to SLAP}
Recent literature, such as the SLAP (Supervised Learning for ASIC Performance) paper, has explored ML for technology mapping. Our approach diverges from SLAP in several key areas to achieve a more lightweight and tightly integrated system:

\begin{enumerate}
    \item \textbf{Model Architecture (MLP vs. GNNs/PairNet):} Frameworks like SLAP often employ complex Graph Neural Networks (GNNs) to capture global circuit topology, followed by heavy PairNet structures to rank cuts. In contrast, our model is a lightweight, 3-layer Multi-Layer Perceptron (MLP). To compensate for the lack of a GNN, we hand-engineered structural features natively computed by ABC (such as \texttt{mffc\_size} and \texttt{slack\_ratio}) to provide the MLP with the necessary topological context.
    \item \textbf{Integration Strategy (Area Blending vs. Hard Filtering):} SLAP typically uses ML as a pre-processing filter—scoring thousands of cuts and discarding the weak ones before passing the subset to the mapper. We took a different approach: \textbf{Inline Area Blending}. Instead of using ML to blindly discard cuts, we score the cuts during ABC's dynamic programming phase and adjust ABC's internal area estimate: $\text{Area}' = \text{Area} \times (1 - \alpha \cdot \tanh(\text{ML\_Score}))$. This acts as a highly intelligent tie-breaker, allowing the ML to gently guide the mapper while preserving ABC's mathematical convergence guarantees.
\end{enumerate}

% ------------------------------------------------------------------
\section{ML-Model Description (ToDo 2)}

The core of the ML-guided mapper is a lightweight Multi-Layer Perceptron (MLP) designed for low-latency inference within the ABC mapping loop.

\subsection{Architecture}
The model is a 3-layer feed-forward neural network:
\begin{itemize}
    \item \textbf{Input Layer}: 9 neurons (corresponding to the selected structural features).
    \item \textbf{Hidden Layer 1}: 64 neurons with ReLU activation and 20\% Dropout for regularization.
    \item \textbf{Hidden Layer 2}: 32 neurons with ReLU activation and 20\% Dropout.
    \item \textbf{Output Layer}: 1 neuron providing a raw scalar score.
\end{itemize}
The model was implemented in PyTorch and subsequently converted to a pure-C implementation for native ABC integration.

\subsection{Feature Engineering \& Selection}
A major component of this project was engineering features that capture cut quality without causing "label leakage." If the model is fed ABC's exact internal heuristics (Area Flow, Delay), it simply learns to copy ABC and provides zero improvement. 

We selected \textbf{9 structural/topological features}:
\texttt{n\_leaves}, \texttt{node\_level}, \texttt{node\_fanout}, \texttt{is\_critical}, \texttt{slack\_ratio}, \texttt{fanout\_adj\_area}, \texttt{area\_per\_leaf}, \texttt{mffc\_size}, and \texttt{mffc\_per\_leaf}. 

Metrics like \texttt{cut\_delay} and exact \texttt{required\_time} were explicitly excluded from the input features. Instead, they were reserved exclusively for constructing the training label.

\subsection{Label Formulation and Training}
Initially, we labeled cuts based on their MFFC size. However, because ABC already optimizes for MFFC, the ML model provided $0.00\%$ QoR improvement over baseline. To force the model to learn a superior objective, we formulated a new label based on the \textbf{Area-Delay Product (ADP)}:
\[ \text{Quality Score} = -(\text{Area Flow} \times \text{Cut Delay}) \]
We used a Pairwise \textbf{MarginRankingLoss} (RankNet). The model is fed pairs of cuts for the same node and is heavily penalized if it incorrectly ranks the cut that has a worse Area-Delay Product.

\subsection{Direct C Codebase Integration}
To avoid the immense I/O overhead of dumping cuts to CSV files and reading them back in python, the trained PyTorch weights were exported to a static C header (\texttt{model\_weights.h}). We wrote a custom C inference engine inside \texttt{ifML.c} to evaluate the MLP using standard floating-point arithmetic. 

% ------------------------------------------------------------------
\section{Methodology: Area Blending Architecture}
Unlike the SLAP framework which often uses ML as a pre-filter, we implemented an \textbf{Inline Area Blending} mechanism. This integrates the ML score directly into ABC's native area-recovery rounds (Mode 2). 

ABC's internal area estimate for a cut is modified as follows:
\[ \text{Area}_{\text{adjusted}} = \text{Area}_{\text{ABC}} \times \left(1 - \alpha \cdot \tanh(\text{ML\_Score})\right) \]
We used a conservative $\alpha = 0.02$ to ensure that the ML acts as an intelligent tie-breaker for structurally similar cuts without disrupting the global convergence of ABC's mapping algorithm.

% ------------------------------------------------------------------
\section{Challenges and Debugging}

Building an ML-integrated mapper presented several severe system-level challenges:

\subsection{The ``Trivial Cut'' Crashing Bug}
During our first integration attempt, we implemented a ``Post-Hoc Override'' architecture where the ML model was allowed to explicitly force ABC to choose its highest-scoring cut at the end of the mapping rounds. 

This led to catastrophic failures and segmentation faults. We discovered the model was assigning extremely high scores to \textbf{trivial cuts} (single-input identity cuts, i.e., \texttt{n\_leaves < 2}). Because these trivial cuts have near-zero delay and area, the ML mathematically deemed them ``perfect.'' When the ML forced ABC to map complex logic cones using identity wires, the mapping graph fractured, causing massive LUT bloat or immediate crashes. 
\textbf{Resolution:} We modified the C code to explicitly ignore and filter out trivial cuts (\texttt{pCut->nLeaves < 2} or \texttt{If\_CutDsd(pCut) == 0}) before ML scoring. 

\subsection{The ``Zero Improvement'' Phase}
Even after fixing the crashes, early iterations of the model yielded exactly 0\% difference from ABC. We traced this to two issues:
1. \textbf{Label Leakage}: The model was trained on MFFC metrics, essentially becoming a clone of ABC's C code. (Resolved by switching to the ADP label).
2. \textbf{Override Timing}: Overriding cuts \textit{after} ABC had converged resulted in broken timing graphs. (Resolved by switching to the \textbf{Area Blending} architecture, blending the ML score directly into \texttt{pCut->Area} during ABC's Mode==2 area-recovery rounds, allowing the ML signal to organically propagate through the graph).

% ------------------------------------------------------------------
\section{Experimental Results}

\subsection{Mapping Time vs. Number of Cuts (ToDo 1)}
To justify the need for ML-guided selection, we measured the baseline CPU mapping time while varying the max-cuts parameter (\texttt{-C}). As shown in Table~\ref{tab:mapping_time}, retaining 16 cuts nearly doubles the runtime compared to the baseline ($C=8$). The goal of our ML integration is to achieve the QoR of high-$C$ mapping without incurring this massive exponential time penalty.

\begin{table}[h!]
\centering
\begin{tabular}{ccc}
\toprule
\textbf{Max Cuts (\texttt{-C})} & \textbf{Total CPU Time (s)} & \textbf{Normalized Time} \\
\midrule
2  & 1.225 & 0.37x \\
4  & 1.855 & 0.55x \\
6  & 2.489 & 0.74x \\
8  & 3.352 & 1.00x (Baseline) \\
12 & 4.666 & 1.39x \\
16 & 6.066 & 1.81x \\
\bottomrule
\end{tabular}
\caption{Aggregate mapping CPU time across 8 benchmarks for varying cut limits.}
\label{tab:mapping_time}
\end{table}

\subsection{Area-Delay QoR Comparison (ToDo 3)}
The final ADP-trained ML model, utilizing a conservative blend factor ($\alpha = 0.02$), was evaluated against the baseline ABC heuristic. Table~\ref{tab:qor} demonstrates that the ML-guided mapper consistently improves area mappings. Across all 20 benchmarks, the ML model achieved an average \textbf{ADP reduction of 2.63\%}, with zero instances of degradation. 

Remarkably, on complex arithmetic designs like the \texttt{square} multiplier, the ML approach reduced the total LUT count by \textbf{12.38\%}, proving that the model successfully learned to identify topologically superior cuts that ABC's localized area-flow heuristics missed.

\begin{table}[h!]
\centering
\resizebox{\textwidth}{!}{
\begin{tabular}{lrrrrrrr}
\toprule
\textbf{Benchmark} & \textbf{ABC LUTs} & \textbf{ML LUTs} & \textbf{Imp \%} & \textbf{ABC ADP} & \textbf{ML ADP} & \textbf{ADP Imp \%} \\
\midrule
adder & 254 & 254 & 0.00 & 12954 & 12954 & 0.00 \\
\textbf{div} & 22030 & 19942 & \textbf{+9.48} & 19033920 & 17229888 & \textbf{+9.48} \\
hyp & 44508 & 41969 & \textbf{+5.70} & 186666552 & 176017986 & \textbf{+5.70} \\
log2 & 8010 & 7778 & \textbf{+2.90} & 616770 & 598906 & \textbf{+2.90} \\
max & 842 & 831 & \textbf{+1.31} & 47152 & 46536 & \textbf{+1.31} \\
multiplier & 5929 & 5683 & \textbf{+4.15} & 314237 & 301199 & \textbf{+4.15} \\
sin & 1464 & 1421 & \textbf{+2.94} & 61488 & 59682 & \textbf{+2.94} \\
\textbf{sqrt} & 5724 & 5180 & \textbf{+9.50} & 5912892 & 5350940 & \textbf{+9.50} \\
\textbf{square} & 3997 & 3502 & \textbf{+12.38} & 199850 & 175100 & \textbf{+12.38} \\
\midrule
mem\_ctrl & 12104 & 11914 & +1.57 & 302600 & 297850 & +1.57 \\
voter & 2826 & 2785 & +1.45 & 48042 & 47345 & +1.45 \\
i2c & 365 & 362 & +0.82 & 1460 & 1448 & +0.82 \\
priority & 219 & 218 & +0.46 & 6789 & 6758 & +0.46 \\
\midrule
\textbf{Arithmetic Mean} & & & \textbf{+2.63\%} & & & \textbf{+2.63\%} \\
\bottomrule
\end{tabular}
}
\caption{QoR Comparison: Baseline ABC vs. ML-Guided ABC. (Benchmarks with 0\% improvement omitted for brevity).}
\label{tab:qor}
\end{table}

\section{Codebase Structure and Implementation Files}
The project is organized into a primary research directory (\texttt{eda\_proj}) containing the core synthesis tool and the machine learning pipeline.

\subsection{Project Directory Organization}
\begin{itemize}
    \item \textbf{abc/}: The modified ABC synthesis tool codebase.
    \item \textbf{benchmarks/}: A collection of AIGER benchmarks (arithmetic and random control).
    \item \textbf{ml\_cut\_project/}: The main machine learning pipeline, scripts, and model assets.
    \begin{itemize}
        \item \textbf{ml/}: Contains trained model weights (\texttt{cut\_model\_mlp.pt}), the feature scaler (\texttt{scaler.pkl}), and training data (\texttt{train\_data.npz}).
        \item \textbf{results/}: Logs and CSV files from QoR comparisons and experiments.
        \item \textbf{abc\_patch/}: Source code patches (\texttt{ifML.c}, \texttt{model\_weights.h}) ready for injection into ABC.
        \item \textbf{data/}: Intermediate raw CSV dumps from ABC's cut profiling.
    \end{itemize}
\end{itemize}

\subsection{Machine Learning Pipeline Scripts (00--09)}
\begin{itemize}
    \item \texttt{00\_setup\_venv.sh}: Configures the Python virtual environment and installs dependencies (PyTorch, Scikit-learn, Pandas).
    \item \texttt{01\_profile\_cuts.sh}: Instruments ABC to dump structural features and mapping quality for millions of cuts across the benchmark suite.
    \item \texttt{02\_analyze\_profile.py}: Provides statistical analysis of the dumped cuts to verify feature distribution and label viability.
    \item \texttt{03\_patch\_and\_build.sh}: Automates the initial patching of the ABC source code and performs a clean build.
    \item \texttt{04\_generate\_training\_data.sh}: Orchestrates the collection of training samples into a unified dataset.
    \item \texttt{05a\_preprocess\_mac.py}: Cleans the data, performs feature engineering (9 structural features), and generates the pairwise training set using the ADP label.
    \item \texttt{05b\_train\_colab.ipynb}: A Google Colab notebook for GPU-accelerated training using RankNet loss.
    \item \texttt{06\_export\_weights\_to\_c.py}: Converts the PyTorch model into a static C header (\texttt{model\_weights.h}) for native execution.
    \item \texttt{07\_install\_ml\_into\_abc.sh}: Injects the generated ML weights and the inference engine into the ABC source tree and recompiles the binary.
    \item \texttt{08\_compare\_qor.py}: Runs the final comparative analysis between baseline ABC and ML-guided ABC.
    \item \texttt{09\_cut\_time\_experiment.py}: Measures the CPU runtime impact of varying cut limits to provide baseline performance metrics.
\end{itemize}

\subsection{Core ABC Modifications}
Rather than relying on file I/O to filter cuts externally, we modified ABC's internal mapper to invoke our ML model natively.
\begin{itemize}
    \item \texttt{src/map/if/ifML.c}: The core inference engine. It contains the \texttt{If\_MLBlendCutArea} function, which executes the forward pass of our MLP and returns the blended area.
    \item \texttt{src/map/if/ifMap.c}: Modified to hook the ML blending function directly into the area recovery mapping rounds (Mode==2).
    \item \texttt{src/map/if/if.h}: Updated with the necessary forward declarations for our custom ML integration.
\end{itemize}

\section{Conclusion}
This project successfully implemented an ML-guided ASIC technology mapper by deeply integrating a lightweight, structural-feature-based MLP directly into ABC's C codebase. By overcoming significant integration hurdles—such as trivial cut crashes and label leakage—and adopting an Area Blending architecture, the model proved capable of outperforming ABC's highly optimized heuristics. The results highlight the potential of machine learning to act not just as an external filter, but as an embedded, mathematical guide for traditional EDA algorithms.

\end{document}
